\chapter{方法设计}
\section{问题定义}

请在此给出研究问题的形式化定义。可使用符号、公式与约束条件说明输入、输出与优化目标。例如，设输入为 $\mathbf{x}$，输出为 $\mathbf{y}$，则目标函数可写为：
\begin{equation}
    \min_{\mathbf{y}} \; \mathcal{L}(\mathbf{x}, \mathbf{y}) + \lambda \Omega(\mathbf{y})
    \label{eq:objective}
\end{equation}
其中，$\mathcal{L}(\cdot)$ 表示主损失项，$\Omega(\cdot)$ 表示正则项，$\lambda$ 为平衡系数。

\section{方法总体框架}

请在此介绍方法总体流程、模块划分与关键设计思路。若需要展示算法流程，可使用算法环境，示例如算法\ref{alg:example}所示。

\begin{algorithm}[htbp]
    \begin{spacing}{1.4}
    \caption{示例算法流程}
    \label{alg:example}
    \algsetup{linenosize=\normalsize}
    \begin{algorithmic}[1]
        \REQUIRE 输入数据 $\mathbf{x}$，超参数 $\alpha$
        \ENSURE 输出结果 $\mathbf{y}$
        \STATE 初始化模型参数 $\theta$
        \FOR{$t = 1$ \TO $T$}
            \STATE 计算损失 $\mathcal{L}(\mathbf{x}, \theta)$
            \STATE 更新参数 $\theta \leftarrow \theta - \alpha \nabla_{\theta}\mathcal{L}$
        \ENDFOR
        \STATE $\mathbf{y} \leftarrow f_{\theta}(\mathbf{x})$
        \STATE \textbf{返回} $\mathbf{y}$
    \end{algorithmic}
    \end{spacing}
\end{algorithm}


\section{关键模块说明}

请在此分小节说明各关键模块的设计细节，包括网络结构、优化策略、损失函数或推理流程等。

\section{实验结果示例}

表\ref{tab:example}给出了一个示例对比结果。请将其替换为你的真实实验表格。

\begin{table}[htbp]
    \centering\small
    \caption{示例对比实验结果（请替换为你的真实数据）}
    \label{tab:example}
    \renewcommand{\arraystretch}{1.2}
    \setlength{\tabcolsep}{9pt}
    \begin{tabular}{lccc}
        \toprule
        \textbf{方法} & \textbf{指标 A}$\downarrow$ & \textbf{指标 B}$\uparrow$ & \textbf{指标 C}$\downarrow$ \\
        \midrule
        基线方法一   & 1.000 & 0.820 & 0.150 \\
        基线方法二   & 0.920 & 0.845 & 0.138 \\
        \midrule
        \textbf{本文方法} & \textbf{0.850} & \textbf{0.880} & \textbf{0.120} \\
        \bottomrule
    \end{tabular}
\end{table}


\section{本章小结}

本章给出了问题定义、方法总体框架与关键模块设计，并通过算法与表格示例展示了模板的常用排版方式。
